\documentclass{article}
\usepackage{amsmath}
\usepackage{graphicx}
\usepackage{biblatex}
\addbibresource{references.bib}

\usepackage{accessibility}
\usepackage{amssymb}

\title{A Rust-Based Implementation and Analysis of the CERA Framework for Climate Model Generalization}
\author{Jules}
\date{\today}

\begin{document}

\maketitle

\begin{abstract}
This paper presents a Rust-based implementation of the CERA (Climate-invariant Encoding through Representation Alignment) framework, a novel machine learning approach designed to improve the generalization of climate models. We detail the architecture of the CERA model, including its autoencoder with 1D convolutional layers and a predictor MLP, and the use of Earth Mover's Distance (EMD) for latent space alignment. The implementation is developed within the \texttt{math_explorer} crate. The end-to-end architecture is verified with tests that confirm the integrity of the data flow and the correctness of component shapes. This work provides a foundation for further research, which would require the integration of a full backpropagation algorithm to enable model training and validation on climate data.
\end{abstract}

\section{Introduction}
Machine learning has shown great promise in advancing climate science, but a major challenge remains in ensuring that models can generalize to out-of-distribution conditions, such as those encountered in a changing climate. The CERA framework, introduced by Liu et al. (2025), offers a solution by learning climate-invariant representations from data. This paper provides a comprehensive Rust implementation of the CERA framework's architecture and validates its structural integrity.

\section{The CERA Framework}
The CERA framework consists of an autoencoder and a predictor that are trained jointly. The key innovation is the use of an Earth Mover's Distance (EMD) loss to align the latent representations of data from different climates \cite{Rubner2000TheRetrieval}.

\subsection{Model Architecture}
The autoencoder uses 1D convolutional layers in its encoder to create a latent representation of the input data, and transposed 1D convolutions in its decoder to reconstruct the input. The predictor is a multi-layer perceptron (MLP) that maps the latent representation to the target variables. This architecture is based on prior work in machine learning for climate science \cite{Yuval2021UsePrecision, Beucler2024Climate-invariantLearning}.

\subsection{Combined Loss Function}
The total loss for the CERA framework is a weighted sum of three components: the reconstruction loss of the autoencoder, the prediction loss of the predictor, and the EMD loss.

\section{Rust Implementation}
The CERA framework is implemented in Rust within the \texttt{math_explorer} crate, in a new module named \texttt{climate}.

\subsection{Project Structure}
The \texttt{climate} module is organized into sub-modules for the CERA model (\texttt{cera.rs}), the autoencoder (\texttt{autoencoder.rs}), the predictor (\texttt{predictor.rs}), the loss functions (\texttt{loss.rs}), and tensor operations (\texttt{tensor_ops.rs}).

\subsection{Core Components}
The implementation includes data structures for the \texttt{Cera}, \texttt{Autoencoder}, and \texttt{Predictor} models. The core logic for data flow, batching, and loss calculation is implemented.

\subsection{Limitations}
A critical component of training a neural network is backpropagation, the process of calculating gradients and updating model weights. The current implementation features a placeholder for the optimization step (\texttt{optimizer_step}) which does not perform true gradient descent. Instead, it simulates weight updates with random noise. This allows for testing the complete data pipeline and architecture, but the model does not learn from the data. A full implementation would require an automatic differentiation engine.

\section{Verification}
The correctness of the implementation's architecture is verified through a series of unit and integration tests. These tests confirm that:
\begin{itemize}
    \item The dimensions of all tensors are handled correctly throughout the forward pass.
    \item The loss functions are calculated correctly based on the model's output.
    \item The data batching and reshaping logic is free of errors.
\end{itemize}
All tests pass, confirming the structural integrity of the implementation and its readiness for the future integration of a backpropagation algorithm.

\section{Conclusion}
This paper has presented a complete architectural implementation of the CERA framework in Rust. We have detailed the model structure, the loss functions, and the data flow, and have verified the implementation with a comprehensive test suite. While the model does not yet learn, this work serves as a robust and validated foundation for future research into climate-invariant machine learning models.

\printbibliography

\end{document}
